\section{Localization and an auxiliary boson--fermion product}
\label{sec:localization}

Localization is a candidate explanation of the cancellations
in Section~\ref{sec:fixed-response}, following Edward Baker's
suggestion. This section gives an exact auxiliary product
and a test for its possible physical origin. It does not
derive the product from a localization complex.

\subsection{The defect-compatible localization question}

Pestun's spherical Wilson-loop calculation organizes
fluctuation cancellations through equivariant kernel and
cokernel multiplicities~\cite[Section~4]{Pestun}.
Baker's open-line discussion instead anticipates reduction
to a one- or two-dimensional defect theory
\cite[Section~3.3]{Baker}. Wang's framework couples a
two-dimensional gauge theory to boundary quantum mechanics
and includes bilocals with endpoint matter and gauge
transport \cite[Section~5.1, equation~(5.5)]{WangDefects}.
The normal-fluctuation determinant is explicitly left
unevaluated in \cite[Section~4.4]{WangDefects}. These results
provide a framework to examine, not the determinant or
boundary response required here.

The conformal open semicircle is an appropriate control.
The constant-Poincare-charge endpoint restriction of
Section~\ref{sec:endpoints} does not exclude a conformal
Killing-spinor construction. Localizing the whole observable
requires its $Q$-closure and boundary conditions compatible
with the bosonic symmetry generated by $Q^2$. Closure of
each factor separately is sufficient, but is not assumed
to be the only possibility.

\subsection{An exact product for the complete prime-free factor}
\label{subsec:bf-product}

Retain one complex bosonic Gaussian at each
$a_n=2n+1/2$, $n=1,\ldots,N$, and one complex Grassmann pair
of weight $p-1/2$. For real $p>1$ define
\begin{align}\label{eq:bf-gaussian}
 Z_N(p)
 &=
 \left(\int\dd\bar\chi\,\dd\chi\,
 e^{-(p-1/2)\bar\chi\chi}\right)
 \prod_{n=1}^N\left(\int_\C\frac{\mathrm d^2z_n}{\pi}
 e^{-(p+a_n)|z_n|^2}\right)\notag\\
 &=(p-1/2)\prod_{n=1}^N(p+a_n)^{-1}.
\end{align}
The Berezin orientation is chosen to give the factor
$p-1/2$. The bosonic integrals converge also at
$p\pm\omega$ for $0\le\omega\le1/2$. One real Gaussian
would give a square root, so statistics and multiplicities
are substantive data.

\begin{proposition}[Auxiliary full-factor identity]
\label{prop:bf-ratio}
For $\Ree p>1$ and $0\le\omega\le1/2$,
\begin{equation}\label{eq:bf-ratio}
 \pi^\omega\lim_{N\to\infty}N^{-\omega}
 \frac{Z_N(p-\omega)}{Z_N(p+\omega)}
 =K_\omega^{<}(p),
\end{equation}
with $K_\omega^{<}$ as in \eqref{eq:small-transfer-factor}.
\end{proposition}
\begin{proof}
Put $z_\pm=(p+1/2\pm\omega)/2$.
The finite product is
\begin{equation}
 \frac{Z_N(p-\omega)}{Z_N(p+\omega)}
 =
 \frac{p-1/2-\omega}{p-1/2+\omega}
 \frac{\Gamma(N+1+z_+)\Gamma(1+z_-)}
      {\Gamma(N+1+z_-)\Gamma(1+z_+)}.
\end{equation}
The first gamma ratio in $N$ is asymptotic to $N^\omega$.
Multiplication by $N^{-\omega}\pi^\omega$ proves the result.
The real Gaussian identity extends analytically to the
stated right half-plane.
\end{proof}

Removing the $n=0$ boson from the earlier gamma model
accounts for the first rational cancellation. The surviving
Grassmann factor supplies the remaining pole ratio.
This is an auxiliary identity with modes prescribed from
the arithmetic target. The coexistence of bosons and
fermions does not specify a supersymmetric action or prove
their pairing by a physical charge. The Berezin factor is
also not a positive Hilbert-space norm.

The normalization in \eqref{eq:bf-ratio} remains external.
In fact, differentiating its finite expression gives
\begin{align}
 -\left.\partial_\omega\log\left[
 \pi^\omega N^{-\omega}
 \frac{Z_N(p-\omega)}{Z_N(p+\omega)}\right]\right|_{\omega=0}
 =\log N-\log\pi+\frac2{p-1/2}
       -2\sum_{n=1}^N\frac1{p+a_n}.
\end{align}
The limit is precisely \eqref{eq:absorbed-symbol}. Reinstating
the removed mode gives the original $w_0$, or equivalently
the fixed finite part of \eqref{eq:contact-finite-part}.
A physical finite-renormalization choice must reproduce it;
counterterm freedom alone does not explain its value.

\subsection{A mode-counting test and the remaining dictionary}
\label{subsec:mode-count}

Differentiate \eqref{eq:absorbed-symbol} in $p$ to remove
an additive $p$-independent generator constant:
\begin{align}\label{eq:mode-count-test}
 \partial_p a_0^{<}(p)
 &=2\sum_{n=1}^\infty\frac1{(p+2n+1/2)^2}
   -\frac2{(p-1/2)^2}\\
 &=2\int_0^\infty u e^{-pu}
 \left[\frac{e^{-5u/2}}{1-e^{-2u}}-e^{u/2}\right]\dd u.
 \notag
\end{align}
Both expressions converge for $\Ree p>1$. They provide a
sharp comparison for a computed virtual mode character,
but do not remove arbitrary $p$-dependent counterterms.
The relative weight $-1/2$ is not by itself a negative
physical energy: the auxiliary integral depends on
$p-1/2$ on a right half-plane. The interpretation of $p$
as a physical boundary or equivariant parameter is missing.

A viable localization calculation must determine the
charge, endpoint conditions, fluctuation complex, boundary
modes and surviving multiplicities, rather than insert
\eqref{eq:mode-count-test}. The reduced theory's correlators
and the observable can contribute alongside its one-loop
determinant; the required factor need not all arise from
that determinant. An explicit preserved charge and the
resulting mode contribution, or a scoped obstruction,
would be a useful first deliverable.

The physical shift also needs its own definition. Under
the localization hypotheses a fixed $Q$-closed expectation
is independent of the coefficient of a $Q$-exact
localization term. That coefficient cannot simply become
$\omega$. A supersymmetric mass, boundary coupling,
admissible contour parameter or source deformation may
vary nontrivially, but its dictionary and all measure and
observable variations must be calculated.

Finally \eqref{eq:bf-ratio} is a partition ratio. A
source-to-response construction must still produce causal
support on arbitrary boundary inputs, the correct identity
limit, and the ordinary physical adjoint. If it retains
the regular readout of \eqref{eq:retarded-readout},
Propositions~\ref{prop:bounded-response} and
\ref{prop:compact-response} still apply.
Signed cohomological cancellation cannot be substituted
for the independent positive cumulative norm. Spatial
gluing, the integer delays beyond $\log2$, and contraction
on a sequence $L_j\to\infty$, $\omega_j\downarrow0$
remain open.

The independently specified hemisphere and scalar-endpoint control is
calculated in Section~\ref{sec:localized-response}. It supplies a
one-sided gamma amplitude, but its fixed protected representation
fails ordinary spatial descent. Its canonical-Hodge alternative
and the distinction between a composite-state character and an
endpoint response are also tested there.
